% 基礎数学 第07回 微分法の基礎（CC BY-SA 4.0）
\documentclass[aspectratio=169]{beamer}
\usepackage{lfs-slides}
\lfscourse{基礎数学}

\title{微分法の基礎}
\subtitle{基礎数学 第07回}
\author{}
\date{}

\begin{document}

\begin{frame}
  \titlepage
\end{frame}

\begin{frame}{今日のねらい}
  \lfsmessage{微分 ＝ その瞬間の「変化の速さ」}
  \bigskip
  \begin{itemize}
    \item 平均変化率から、瞬間の変化率（微分係数）へ
    \item $x^n$ の微分の公式を使えるようになる
    \item 接線の方程式を求める
    \item 工学では：位置を微分すると速度になる
  \end{itemize}
\end{frame}

\section{平均変化率}

\begin{frame}{平均変化率：「平均してどれだけ増えたか」}
  \begin{columns}[c]
    \column{0.5\textwidth}
      \begin{center}
        % 平均変化率（割線）の図：y = x^2 の x=1 から x=3 まで
\begin{tikzpicture}[xscale=1.0,yscale=0.5]
  \draw[->] (-0.5,0) -- (4,0) node[right] {$x$};
  \draw[->] (0,-0.5) -- (0,10) node[above] {$y$};
  \draw[thick,blue,domain=0:3.15,samples=60] plot (\x,{\x*\x}) node[right] {$y=x^2$};
  % 割線：(1,1) と (3,9) を通る（傾き 4）
  \draw[red,thick] (0.5,-1) -- (3.4,10.6);
  \fill (1,1) circle (2.5pt) node[left] {$(1,\,1)$};
  \fill (3,9) circle (2.5pt) node[left] {$(3,\,9)$};
  \draw[dashed] (1,1) -- (3,1) -- (3,9);
  \node[below] at (2,1) {$x$ の増加 $2$};
  \node[right] at (3,5) {$y$ の増加 $8$};
\end{tikzpicture}

      \end{center}
    \column{0.5\textwidth}
      $x$ が $a$ から $b$ まで変わるときの
      \[
        \text{平均変化率} = \frac{f(b) - f(a)}{b - a}
      \]
      \medskip
      $f(x) = x^2$、$1 \to 3$ なら
      \[
        \frac{9 - 1}{3 - 1} = \frac{8}{2} = 4
      \]
      \medskip
      {\small 図の赤い直線（割線）の傾きと同じ}
  \end{columns}
\end{frame}

\begin{frame}{例：ボールが転がった距離}
  \lfsmessage{平均変化率 ＝ 平均の速さ}
  \medskip
  時刻 $t$ 秒までに $s(t) = 5t^2$ メートル進むとき、$t = 1$ から $t = 3$ までの平均の速さは
  \[
    \frac{s(3) - s(1)}{3 - 1} = \frac{45 - 5}{2} = 20 \ \text{(m/s)}
  \]
  \begin{alertblock}{でも知りたいのは…}
    「$t = 1$ の\textbf{その瞬間}の速さ」は？ \quad → 区間をどんどん短くする
  \end{alertblock}
\end{frame}

\section{微分係数}

\begin{frame}{区間を縮めていく}
  \begin{columns}[c]
    \column{0.52\textwidth}
      \begin{center}
        % 割線が接線に近づく図：y = x^2 の点 (1,1) で、相手の点 (1+h, (1+h)^2) を近づける
\begin{tikzpicture}[xscale=1.35,yscale=0.6]
  \begin{scope}
    \clip (-0.4,-0.8) rectangle (3.2,6.6);
    \draw[thick,blue,domain=0:2.55,samples=60] plot (\x,{\x*\x});
    % 割線（h = 1.5, 1, 0.5）：(1,1) と (1+h,(1+h)^2) を通る。傾き 2 + h
    \foreach \h/\c in {1.5/red!30, 1/red!55, 0.5/red!85} {
      \draw[\c,thick,domain=0.6:{1+\h}] plot (\x,{(2+\h)*(\x-1)+1});
      \fill[\c] ({1+\h},{(1+\h)*(1+\h)}) circle (2.5pt);
    }
    % 接線（傾き 2）
    \draw[black,very thick,domain=0.3:3.0] plot (\x,{2*(\x-1)+1});
  \end{scope}
  \draw[->] (-0.4,0) -- (3.3,0) node[right] {$x$};
  \draw[->] (0,-0.8) -- (0,6.8) node[above] {$y$};
  \fill (1,1) circle (2.5pt) node[below right] {$(1,\,1)$};
  \node[right] at (3.0,5.0) {接線};
  \foreach \h in {1.5, 1, 0.5} {
    \node[above left,inner sep=1pt] at ({1+\h},{(1+\h)*(1+\h)}) {\small $h=\h$};
  }
\end{tikzpicture}

      \end{center}
    \column{0.48\textwidth}
      $f(x) = x^2$、$x = 1$ から $1 + h$ まで
      \medskip
      \begin{center}
        \begin{tabular}{rr}
          \toprule
          $h$ & 平均変化率 \\
          \midrule
          $1$ & $3$ \\
          $0.1$ & $2.1$ \\
          $0.01$ & $2.01$ \\
          $0.001$ & $2.001$ \\
          \bottomrule
        \end{tabular}
      \end{center}
      \medskip
      $h \to 0$ で \textbf{$2$ に近づく}
  \end{columns}
\end{frame}

\begin{frame}{微分係数の定義}
  \lfsmessage{$\displaystyle f'(a) = \lim_{h \to 0} \frac{f(a+h) - f(a)}{h}$}
  \bigskip
  \begin{itemize}
    \item $f'(a)$ を「$x = a$ における\textbf{微分係数}」とよぶ
    \item $\lim\limits_{h \to 0}$ は「$h$ を $0$ に限りなく近づけたときの値」
    \item 図では：割線が\textbf{接線}に近づく → $f'(a)$ は接線の傾き
  \end{itemize}
\end{frame}

\begin{frame}{定義どおりに計算してみる}
  $f(x) = x^2$ の $x = a$ における微分係数
  \begin{align*}
    \frac{f(a+h) - f(a)}{h}
      &= \frac{(a+h)^2 - a^2}{h}
       = \frac{2ah + h^2}{h} \\
      &= 2a + h
      \quad \xrightarrow{\ h \to 0\ } \quad 2a
  \end{align*}
  \begin{block}{結果}
    $f'(a) = 2a$ \qquad （$a = 1$ なら $2$：表の結果と一致）
  \end{block}
\end{frame}

\section{導関数と公式}

\begin{frame}{導関数}
  \lfsmessage{$a$ を $x$ にしたもの ＝ 導関数 $f'(x)$}
  \medskip
  \begin{itemize}
    \item $f(x) = x^2$ の導関数は $f'(x) = 2x$
    \item 導関数を求めることを「\textbf{微分する}」という
    \item 書き方はいろいろ：$f'(x)$，\ $y'$，\ $\dfrac{dy}{dx}$
  \end{itemize}
\end{frame}

\begin{frame}{微分の公式}
  \begin{columns}[t]
    \column{0.5\textwidth}
      \begin{block}{べき関数}
        \[
          (x^n)' = n\,x^{n-1}
        \]
        {\small 例：$(x^3)' = 3x^2$，\ $(x)' = 1$}
      \end{block}
      \begin{block}{定数}
        \[
          (c)' = 0
        \]
      \end{block}
    \column{0.5\textwidth}
      \begin{block}{定数倍と和}
        \[
          \{k f(x)\}' = k f'(x)
        \]
        \[
          \{f(x) + g(x)\}' = f'(x) + g'(x)
        \]
      \end{block}
      \begin{block}{三角関数（参考）}
        $(\sin x)' = \cos x$，\ $(\cos x)' = -\sin x$\\
        {\small （$x$ は弧度法）}
      \end{block}
  \end{columns}
\end{frame}

\begin{frame}{公式を使う}
  \[
    (3x^2 - 4x + 1)' = 3 \cdot 2x - 4 \cdot 1 + 0 = 6x - 4
  \]
  \bigskip
  \[
    (x^3 - 2x^2 + 5x)' = 3x^2 - 4x + 5
  \]
  \bigskip
  \begin{block}{手順}
    項ごとに分けて、\ 係数はそのまま、\ $x^n \to n x^{n-1}$
  \end{block}
\end{frame}

\section{接線}

\begin{frame}{接線の方程式}
  \lfsmessage{点 $(a,\,f(a))$ を通り、傾き $f'(a)$ の直線}
  \[
    y - f(a) = f'(a)\,(x - a)
  \]
  \medskip
  例：$y = x^2$ の点 $(1,\,1)$ における接線
  \begin{itemize}
    \item $f'(x) = 2x$ より $f'(1) = 2$
    \item $y - 1 = 2(x - 1)$ \ → \ $y = 2x - 1$
  \end{itemize}
\end{frame}

\section{工学では}

\begin{frame}{位置を微分すると速度}
  \lfsmessage{速度 ＝ 位置の時間変化率}
  \medskip
  ボールの位置 $s(t) = 5t^2$（m）のとき
  \[
    v(t) = s'(t) = 10t \ \text{(m/s)}
  \]
  $t = 1$ の瞬間の速さは $10$ m/s、\ $t = 3$ では $30$ m/s
  \begin{block}{使われる場面}
    自動車・ロボットの制御（速さ・加速度の調整）、\
    物理シミュレーション、\
    センサーの値の変化から異常を見つける など
  \end{block}
\end{frame}

\begin{frame}{まとめ}
  \begin{itemize}
    \item 平均変化率 $\dfrac{f(b) - f(a)}{b - a}$ ＝ 割線の傾き
    \item 微分係数 $f'(a) = \lim\limits_{h \to 0} \dfrac{f(a+h) - f(a)}{h}$ ＝ 接線の傾き
    \item 公式：$(x^n)' = n x^{n-1}$、定数は $0$、項ごとに微分できる
    \item 接線：$y - f(a) = f'(a)(x - a)$
    \item 位置を微分すると速度
  \end{itemize}
  \bigskip
  \lfsmessage{次回：積分法の基礎（微分の逆）}
\end{frame}

\end{document}
